%% kapa: exam=1 round=37 year=2026 period=1 form=A subject=real_estate subject_ko=부동산학원론 seq=31 number=111
%% source: Q-Net 감정평가사 기출문제 (cst00309) artlSeq=5258924 file=raw/1008/5258924/2026년 제37회 감정평가사 1차 1교시.pdf
%% conversion: kordoc 4.24.51 -> tools/gen1.py (auto); review=proofread-sonnet; flags=none
\begin{question}{111}
부동산관리방식에 관한 설명으로 옳지 않은 것은?
\begin{choices}
  \item 자기관리는 관리의 전문성이 부족하나 비용이 적게 들고 임대료나 관리비의 객관적 산정이 용이하다.
  \item 위탁관리는 전문관리업자에게 위탁하여 관리하는 방식으로 대형건물의 관리에 유용하다.
  \item 혼합관리는 필요한 부분만 위탁하여 관리비 절감 및 효과적 운용이 가능하나 책임소재가 불분명하다.
  \item 자기관리는 고객에 대한 정보관리가 쉽고 부동산에 관한 기밀유지와 보안관리가 용이하다.
  \item 위탁관리는 부동산 환경의 변화에 능동적으로 대응할 수 있으나 하자발생시 신속한 대응이 어렵다.
\end{choices}
\end{question}
\answer{1}
